\documentclass[conference]{IEEEtran}

\usepackage{amsmath,amssymb}
\usepackage{booktabs}
\usepackage{graphicx}
\usepackage{hyperref}
\usepackage{microtype}
\usepackage{xcolor}

\newcommand{\method}{DynNav-H}
\newcommand{\Rret}{R_{\mathrm{return}}}
\newcommand{\Aset}{\mathcal{A}}
\newcommand{\Ssafe}{\mathcal{S}_{\mathrm{safe}}}
\newcommand{\todo}[1]{\textcolor{red}{[TODO: #1]}}

\title{When the Same Place Is Not the Same State:\\History-Conditioned Safe-Return Planning under Action-Triggered Topology Hazards}

\author{\IEEEauthorblockN{Anonymous Author(s)}
\IEEEauthorblockA{Affiliation withheld for review}}

\begin{document}
\maketitle

\begin{abstract}
A mobile robot can reach the same geometric location through two paths while retaining different future return options if its earlier actions alter the probability of later environmental closures. A state-only navigation model cannot represent this distinction. We study action-triggered topology hazards: traversing selected directed transitions activates stochastic future closures of cells that may be critical for returning to a designated safe set. We formulate safe-return reliability on an augmented state consisting of robot position and activated hazards, give an exact small-world oracle under independent Bernoulli closures, and plan with A* using a soft recoverability cost. A counterfactual benchmark shows that histories ending at the same cell can differ in return reliability by as much as 0.9 while a state-only marginal-risk model necessarily assigns them the same value. In an analytic route-switch family, the exact planner matches all 75 predicted switch decisions. In paired stochastic execution over 1000 seeds per condition, the shortest and state-only marginal-risk planners exhibit identical irreversible-failure rates of 0.512, 0.881, and 0.992 at closure probabilities 0.2, 0.5, and 0.8, whereas the history-aware and hard safe-return planners avoid all trigger activations in this controlled family. A critical-cut approximation accelerates exact return evaluation by about 177$\times$ at 12 series hazards, but we also construct a joint-cut counterexample in which this approximation is intentionally optimistic. The results isolate a specific information advantage of history conditioning while making explicit the limits of the current discrete, known-model evaluation.
\end{abstract}

\section{Introduction}
Planning for collision avoidance or low expected traversal risk is not equivalent to preserving the ability to recover. A robot may enter a corridor safely under the current map and still make a commitment that leaves no return route after a later closure. Recent work has formalized safe-return policies, contingency feasibility, returnability, backup-plan safety, and predictive commitment in uncertain environments \cite{guo2023safereturn,stephens2024safeexploration,jung2025contingency,srirangam2026scramppi,tao2024backup,han2026rcsp}. These results rule out any broad novelty claim for ``planning while keeping a way back.''

This paper asks a narrower question: \emph{what if the robot's own path history changes which future environmental topology hazards are active?} Consider two trajectories that end at the same geometric cell. One traversed a transition that causes a door, gate, or passage behind the robot to become likely to close; the other avoided that transition. If the closure process is action-triggered, the two trajectories have different future return-connectivity despite sharing the same current position and static map. The relevant planning state is therefore not position alone.

Decision-dependent uncertainty is a mature topic beyond robotics \cite{doan2022endogenous}, and robot planning has considered utilization-driven component degradation \cite{baldes2024fallible}, dynamic door-closing workspaces \cite{liu2021door}, and nondeterministic environmental events \cite{perret1996events}. Our contribution is not the generic idea that actions influence future uncertainty. Instead, we isolate and experimentally test the information lost when an action-triggered stochastic process changes the \emph{environment's future return-connectivity} but a planner retains only state-only marginal closure probabilities.

The contributions are:
\begin{itemize}
  \item a compact model of action-triggered topology hazards and an augmented planning state $(x,\Aset)$ that retains which future closures have been activated;
  \item an exact safe-return reliability oracle for independent future closures and an A* objective that trades path length against history-conditioned return fragility;
  \item a counterfactual same-state/different-history test showing an information gap that a state-only marginal model cannot represent;
  \item analytic, procedural, stochastic, scaling, held-out, and adversarial benchmarks against shortest-path, state-only probability-aware, critical-cut, and hard safe-return baselines;
  \item an explicit failure analysis showing when the scalable critical-cut approximation becomes optimistic under joint multi-cell cutsets.
\end{itemize}

\section{Related Work}
\subsection{Safe return and contingency feasibility}
Guo \emph{et al.} require the existence of a policy that returns an MDP to safe states with high probability and jointly synthesize outbound and return behavior \cite{guo2023safereturn}. Safe exploration methods likewise evaluate reachability and returnability probabilities before visiting uncertain regions \cite{stephens2024safeexploration}. Backup Plan Constrained MPC optimizes primary and alternative mission feasibility \cite{tao2024backup}. Contingency-MPPI embeds contingency generation inside nominal planning, and SCRAMPPI uses Hamilton--Jacobi reachability to certify that a safe-set trajectory exists from every state on the nominal plan \cite{jung2025contingency,srirangam2026scramppi}. These methods are stronger than our grid abstraction as general safety or control frameworks. Our distinction is representational: the environment hazard set itself depends on the executed transition history.

\subsection{Predictive commitment and dynamic topology}
RCSP studies commands that are locally safe but commit a robot to a bottleneck that moving obstacles may soon close \cite{han2026rcsp}. Dynamic door-closing work has used POMDP belief-tree reasoning for self-protective manipulation behavior \cite{liu2021door}, while nondeterministic planning has long represented exogenous events and robust plans \cite{perret1996events}. These works motivate proactive reasoning about future changes; they do not by themselves establish that our action-triggered return-topology abstraction is novel.

\subsection{Endogenous uncertainty and degradation}
Optimization under endogenous uncertainty explicitly models probability distributions that depend on decisions \cite{doan2022endogenous}. Baldes \emph{et al.} model utilization-driven, state-action-dependent actuator failures and their long-term effect on future control capability \cite{baldes2024fallible}. Consequently, neither ``history matters'' nor ``actions can reduce future capabilities'' is a contribution here. We instead focus on a falsifiable robotics question: whether history contains return-connectivity information that is discarded when environmental closure probabilities are represented only as state-only marginals.

\section{Problem Formulation}
Let $G=(V,E)$ be a traversable grid graph, $x_t\in V$ the robot position, and $\Ssafe\subseteq V$ a safe set. A topology hazard
\begin{equation}
h_i=(e_i,c_i,p_i)
\end{equation}
contains a directed trigger transition $e_i\in E$, a currently traversable closure cell $c_i\in V$, and a future closure probability $p_i\in[0,1]$. Traversing $e_i$ activates hazard $i$. For path history $H_t=(x_0,\ldots,x_t)$, define
\begin{equation}
\Aset(H_t)=\{i\mid e_i\text{ occurs in }H_t\}.
\end{equation}
The augmented planning state is $z_t=(x_t,\Aset(H_t))$.

For a fixed augmented state, each active closure cell is sampled independently as a Bernoulli event. The current robot cell is conditioned traversable at the evaluation instant. Let $C$ denote the random set of realized closures. We define exact safe-return reliability
\begin{equation}
\Rret(x,\Aset)=
\Pr\!\left[\exists\text{ path in }G\setminus C\text{ from }x\text{ to }\Ssafe\right].
\label{eq:return}
\end{equation}
For small active sets, Eq.~\eqref{eq:return} is computed exactly by enumerating closure realizations and testing connectivity.

\subsection{Why geometric state alone is insufficient}
A state-only marginal model maps position $x$ and a fixed marginal hazard field to a single return estimate $\widetilde R(x)$. It therefore cannot assign distinct values to two histories that end at the same $x$.

\textbf{Proposition 1 (history information gap).}
There exists a grid, safe set, trigger hazard, and histories $H^a,H^b$ with the same endpoint $x$ such that
\begin{equation}
\Rret(x,\Aset(H^a))\neq\Rret(x,\Aset(H^b)),
\end{equation}
while any deterministic state-only marginal model returns the same $\widetilde R(x)$ for both histories.

\emph{Proof.} Let one history traverse a trigger that activates closure of the unique bridge to the safe set with probability $p>0$, and let the other history reach the same endpoint without traversing that trigger. Then the triggered history has return reliability $1-p$ and the non-triggered history has reliability $1$. Since both end at the same $x$, a state-only function of $x$ and fixed marginals must return one common value. \hfill$\square$

\section{Planning Methods}
\subsection{Exact history-aware A*}
Search is performed over augmented states $(x,\Aset)$. For transition $z\rightarrow z'$, the soft-history objective uses
\begin{equation}
c(z,z')=1+\lambda\left(1-\Rret(z')\right),
\label{eq:cost}
\end{equation}
with nonnegative recoverability weight $\lambda$. The online return-oracle calls used to select the route are included in measured planning latency; post-hoc diagnostics are excluded.

\subsection{State-only probability-aware baseline}
The state-only baseline receives the same closure cells and marginal probabilities $p_i$ but drops trigger history. It plans using return-path reliability under this fixed marginal field. Thus it is not an uninformed shortest-path baseline: it knows where future closures may occur and with what marginal probability, but it cannot know whether a particular path actually activates them.

\subsection{Hard safe-return baseline}
A strong constraint baseline accepts a successor only when
\begin{equation}
\Rret(z')\geq\tau,
\end{equation}
for a threshold $\tau$. This abstracts safe-return and contingency methods that enforce backup feasibility rather than optimize a soft penalty.

\subsection{Critical-cut approximation}
Exact enumeration scales exponentially in the number of active hazards. We therefore evaluate a structural approximation that removes each active hazard cell individually and identifies those whose single removal disconnects the robot from $\Ssafe$. Multiplying the survival probabilities of these individually critical cells gives
\begin{equation}
R_{\mathrm{cut}}=\prod_{i\in K}(1-p_i).
\end{equation}
Because multi-cell cutsets may disconnect the graph even when no individual cell is critical, $R_{\mathrm{cut}}$ is an upper bound in general. We treat this approximation as scalability infrastructure, not as a novel graph-theoretic contribution.

\section{Experimental Design}
All stochastic comparisons use common random numbers: a seed generates one latent draw for every declared closure event, and a planner realizes only the events it activated. This preserves paired event-index outcomes even when planners activate different subsets. Irreversible failure is operationalized as post-closure inability to reach the safe set in the resulting grid. Failed trials remain in binary denominators; undefined continuous metrics are counted explicitly rather than silently removed.

We evaluate six complementary questions.
\begin{enumerate}
  \item \emph{Information gap:} can identical geometric endpoints have different exact return reliability solely because of history?
  \item \emph{Analytic route switching:} does the planner follow a predicted length--fragility boundary rather than a post-hoc tuned pattern?
  \item \emph{Stochastic execution:} does history conditioning change irreversible-failure outcomes under sampled future closures?
  \item \emph{State-only falsification:} does a probability-aware planner with the same marginals reproduce the history-aware behavior?
  \item \emph{Scaling:} what computation is introduced by exact history-conditioned evaluation, and what does the cut approximation save?
  \item \emph{Approximation boundary:} can a joint cutset produce a route-choice failure for the cut approximation?
\end{enumerate}

\section{Results}
\subsection{Same state, different history}
The counterfactual benchmark sweeps $p\in\{0.1,0.3,0.5,0.7,0.9\}$. The safe history has exact return reliability $1$, while the trigger-taking history has reliability $1-p$. A state-only marginal model assigns $1-p$ to both because it cannot condition on whether the trigger was traversed. Across the five cases, the mean history separation is $0.5$ and the maximum is $0.9$; the triggered-history invariant is matched with zero numerical error.

\subsection{Analytic phase boundary}
In a controlled world with a direct path of length $4$ and a trigger-avoiding detour of length $4+2d$, exactly one direct transition incurs fragility $\lambda p$. The predicted route switch is therefore
\begin{equation}
\lambda p>2d.
\end{equation}
Across 75 combinations of detour depth, closure probability, and weight, the exact history-aware planner matches the analytic route choice in all 75 cases. The shortest-path baseline never detours.

\subsection{Stochastic execution and state-only baseline}
Table~\ref{tab:execution} reports 1000 paired execution seeds at each closure probability for a three-module series family. The shortest and state-only probability-aware planners select the same trigger-taking path and consequently have identical binary outcomes. The exact-history, critical-cut, and hard $\tau=0.9$ planners avoid all triggers for the high recoverability weight used in this controlled experiment.

\begin{table}[t]
\centering
\caption{Controlled stochastic execution, 1000 paired seeds per $p$. ``Fail'' is post-closure return infeasibility.}
\label{tab:execution}
\begin{tabular}{lccc}
\toprule
Planner & $p=0.2$ & $p=0.5$ & $p=0.8$\\
\midrule
Shortest & 0.512 & 0.881 & 0.992\\
State-only marginal & 0.512 & 0.881 & 0.992\\
History exact & 0 & 0 & 0\\
History cut & 0 & 0 & 0\\
Hard return $\tau=.9$ & 0 & 0 & 0\\
\bottomrule
\end{tabular}
\end{table}

For the history-exact comparison against shortest, paired risk differences are $-0.512$, $-0.881$, and $-0.992$. Exact McNemar tests are extremely small because all discordant outcomes favor history-aware planning in this controlled family. In contrast, the state-only marginal baseline has zero discordant pairs and McNemar $p=1$ at every probability. This null comparison is central: marginal probability awareness alone does not explain the history-aware route choice.

\subsection{Approximation and scaling}
On a series-critical family, the critical-cut estimate is exact by construction. At 12 hazards, exact enumeration takes 61.8~ms in the retained artifact while the cut estimate takes 0.348~ms, a speedup of approximately 177$\times$. This result must not be generalized beyond the family on which the cut structure is exact.

Search-level timing after excluding post-hoc diagnostic profiling is shown in Table~\ref{tab:scaling}. At six repeated commitment modules, shortest augmented search takes 0.119~ms, history-cut search 3.93~ms, and exact-history search 8.99~ms. The measurement includes online return queries that affect route choice.

\begin{table}[t]
\centering
\caption{Augmented-state search at six modules.}
\label{tab:scaling}
\begin{tabular}{lrrr}
\toprule
Planner & Nodes & Path length & Planning ms\\
\midrule
Shortest augmented & 25 & 24 & 0.119\\
History cut & 72 & 36 & 3.93\\
History exact & 72 & 36 & 8.99\\
\bottomrule
\end{tabular}
\end{table}

\subsection{Held-out heterogeneous probabilities}
\todo{Replace this paragraph with CI-verified PR \#144 held-out artifact values. The protocol is frozen before inspection: 6--8 modules with per-trigger probability vectors not used in the development sweeps, 500 paired seeds per scenario, and the same five planners.}

\subsection{Joint-cut adversarial case}
\todo{Replace this paragraph with CI-verified PR \#144 results. The adversarial topology contains two parallel return corridors; no single hazard is critical, but simultaneous closure of both disconnects the safe set. The expected result is that exact history-aware planning detours for sufficiently large $\lambda p^2$, whereas the individually-critical cut approximation can remain maximally optimistic.}

\section{Discussion}
The experiments support a narrow claim. When actions activate future topology hazards, path history can contain information about safe-return connectivity that is absent from a state-only marginal hazard map. This is not equivalent to claiming that history-aware planning is universally safer, that soft recoverability costs dominate hard contingency constraints, or that the proposed model subsumes general dynamic navigation.

The state-only null result is more informative than comparison only against shortest path. Both planners know the same closure locations and marginal probabilities. Their divergence arises because one conditions on the executed trigger set and the other does not. Conversely, the hard safe-return baseline matches the exact-history planner in the high-penalty series experiment, showing that our current evidence establishes an information mechanism rather than superiority over all safe-return formulations.

The exact oracle is useful for small controlled worlds and for validating approximations, but its exponential closure enumeration is not a scalable navigation solution. The critical-cut approximation removes most of this cost in series-articulation families. Its joint-cut failure is therefore an important design boundary: fast structural approximations require either multi-cell cut reasoning, scenario sampling, or another estimator when redundant return paths share correlated failure modes.

\section{Limitations}
The current study uses discrete grid connectivity, known trigger locations, known closure probabilities, and mostly independent Bernoulli future closures. A scenario-based oracle demonstrates that equal marginals can be optimistic under correlated closures, but correlated uncertainty is not yet integrated into the online planner. The stochastic execution experiments sample topology events after commitment; they do not model continuous robot dynamics, perception delay, collision, or contact.

The repository contains a ROS~2 Jazzy/Nav2/Gazebo dynamic route-invalidation protocol, but the history-conditioned planner evaluated in this paper has not yet been integrated into that execution stack. We therefore make no hardware, collision-avoidance, kinodynamic-safety, or formal runtime-guarantee claim. The result should be interpreted as algorithmic and mechanism evidence for a specific representational gap.

Finally, literature search cannot prove novelty. Safe return, contingency planning, predictive commitment, action reversibility, endogenous uncertainty, dynamic doors, and utilization-driven degradation all have substantial prior art \cite{guo2023safereturn,jung2025contingency,srirangam2026scramppi,han2026rcsp,baldes2024fallible,liu2021door,doan2022endogenous}. The defensible contribution is the explicit action-triggered environmental return-topology formulation and its state-only information-gap evaluation, subject to further prior-art review.

\section{Conclusion}
We studied navigation in which the robot's own transition history activates stochastic future closures that can destroy return connectivity. In this setting, the same geometric cell can correspond to different recoverability because the environment's future topology depends on how the robot arrived. Augmenting search with the activated hazard set captures this distinction; a state-only marginal-risk baseline cannot. Exact planning follows an analytic route-switch boundary and eliminates failures in a controlled stochastic series family, while a scalable critical-cut approximation offers substantial speedup but fails on joint cutsets. These results establish a reproducible mechanism-level case for history-conditioned safe-return reasoning and clearly delimit the additional evidence required for broader robotics claims.

\bibliographystyle{IEEEtran}
\bibliography{references}

\end{document}
